\chapter{Security, safety and ethics}\label{ch:security}

\section{Emissions are the main risk}

On a modern battlefield, a transmission is a beacon. Near is designed to need few, short transmissions:
\begin{itemize}
  \item one report, then a few questions, each chosen to carry the most information. Every question costs a
        transmission, so we ask the fewest;
  \item measurement packets of about 30 bytes, about 2 s of tones, instead of seconds of raw audio;
  \item a phone that guides offline from one route packet and only transmits at turns or on arrival. The earlier
        connected design polled the server every 3 s, which in the field would be a beacon. It's kept for
        connected use only.
\end{itemize}
In radio mode, Near adds no emitter of its own. The radio the unit already uses carries everything.

\section{Confidentiality and integrity}

\begin{itemize}
  \item \textbf{Nothing leaves the laptop.} Speech-to-text, language model, maps, sound models and voices all run locally. No
        cloud service sees the audio or the positions.
  \item \textbf{Packets are not encrypted yet.} Military radios encrypt the voice channel, and the tones ride inside it.
        On an unencrypted radio, a listener could decode a packet. Authenticated encryption with a pre-shared unit key
        (e.g. AES-CCM, 4--8 bytes of tag) fits the packet budget and is on the roadmap.
  \item \textbf{Spoofing.} A hostile station could send fake reports or packets. Callsign checks, sequence numbers (already in the
        packet), and later authentication tags address this. The coordinator approves every spoken reply.
  \item \textbf{GPS spoofing.} A live GPS fix is just one constraint among many, with an outlier floor. If it disagrees with the
        landmarks, it loses.
\end{itemize}

\section{Medical safety}

Near organises what the field reports. It doesn't diagnose, prescribe or set priorities on its own authority:
\begin{itemize}
  \item every medical rule is a visible, configurable placeholder for medical validation;
  \item ``no data'' is never shown as ``fine'': unknown body parts are orange, unknown MARCH items grey;
  \item inferred values (mobility, posture) are marked as inferred, and conflicts are flagged, not resolved silently;
  \item voice stress changes how questions are asked, never the casualty's priority;
  \item a heart rate is shown only from a chest recording, never estimated from the voice.
\end{itemize}
Depending on use, decision-support software can fall under medical-device rules (in the EU, the MDR). That
classification must be settled before any clinical use.

\section{Dual use and export control}

Near locates and evacuates casualties and maps threats reported by friendly units. It doesn't target anything.
Friendlies only: the design never locates other people, and danger zones mark areas to avoid. Software for military use
may still fall under export-control rules, and any distribution outside the team needs that check.

\section{Failure behaviour}

\begin{itemize}
  \item LLM down $\rightarrow$ rules alone. Sound models missing $\rightarrow$ voice-only positioning. LiDAR missing $\rightarrow$
        building heights and woods.
  \item No position possible $\rightarrow$ ``no position'', never a guess. A town outside the mapped area is detected when named, and
        no position is invented.
  \item Ambiguous position $\rightarrow$ all candidate areas shown and spoken, and no medic route until it's resolved.
\end{itemize}
